% Options for packages loaded elsewhere
\PassOptionsToPackage{unicode}{hyperref}
\PassOptionsToPackage{hyphens}{url}
\PassOptionsToPackage{dvipsnames,svgnames,x11names}{xcolor}
\documentclass[
  10pt,
  a4paper,
]{article}
\usepackage{xcolor}
\usepackage[margin=18mm]{geometry}
\usepackage{amsmath,amssymb}
\setcounter{secnumdepth}{-\maxdimen} % remove section numbering
\usepackage{iftex}
\ifPDFTeX
  \usepackage[T1]{fontenc}
  \usepackage[utf8]{inputenc}
  \usepackage{textcomp} % provide euro and other symbols
\else % if luatex or xetex
  \usepackage{unicode-math} % this also loads fontspec
  \defaultfontfeatures{Scale=MatchLowercase}
  \defaultfontfeatures[\rmfamily]{Ligatures=TeX,Scale=1}
\fi
\usepackage{lmodern}
\ifPDFTeX\else
  % xetex/luatex font selection
  \setmainfont[]{DejaVu Serif}
  \setsansfont[]{DejaVu Sans}
  \setmonofont[]{DejaVu Sans Mono}
\fi
% Use upquote if available, for straight quotes in verbatim environments
\IfFileExists{upquote.sty}{\usepackage{upquote}}{}
\IfFileExists{microtype.sty}{% use microtype if available
  \usepackage[]{microtype}
  \UseMicrotypeSet[protrusion]{basicmath} % disable protrusion for tt fonts
}{}
\makeatletter
\@ifundefined{KOMAClassName}{% if non-KOMA class
  \IfFileExists{parskip.sty}{%
    \usepackage{parskip}
  }{% else
    \setlength{\parindent}{0pt}
    \setlength{\parskip}{6pt plus 2pt minus 1pt}}
}{% if KOMA class
  \KOMAoptions{parskip=half}}
\makeatother
\usepackage{longtable,booktabs,array}
\newcounter{none} % for unnumbered tables
\usepackage{calc} % for calculating minipage widths
% Correct order of tables after \paragraph or \subparagraph
\usepackage{etoolbox}
\makeatletter
\patchcmd\longtable{\par}{\if@noskipsec\mbox{}\fi\par}{}{}
\makeatother
% Allow footnotes in longtable head/foot
\IfFileExists{footnotehyper.sty}{\usepackage{footnotehyper}}{\usepackage{footnote}}
\makesavenoteenv{longtable}
\setlength{\emergencystretch}{3em} % prevent overfull lines
\providecommand{\tightlist}{%
  \setlength{\itemsep}{0pt}\setlength{\parskip}{0pt}}
\usepackage{xeCJK}
\setCJKmainfont[AutoFakeBold=1.4]{Droid Sans Fallback}
\setCJKsansfont[AutoFakeBold=1.4]{Droid Sans Fallback}
\setCJKmonofont{Droid Sans Fallback}
\emergencystretch=3em
\setlength{\parindent}{0pt}
\xeCJKDeclareCharClass{Default}{"2013,"00B7}
\usepackage{newunicodechar}
\newunicodechar{ℱ}{\ensuremath{\mathcal{F}}}
\renewcommand{\contentsname}{目录}

\usepackage{xurl}
\usepackage{bookmark}
\IfFileExists{xurl.sty}{\usepackage{xurl}}{} % add URL line breaks if available
\urlstyle{same}
\hypersetup{
  pdftitle={He--Kruczenski 有限原型：底层证明与适用范围},
  colorlinks=true,
  linkcolor={Maroon},
  filecolor={Maroon},
  citecolor={Blue},
  urlcolor={Blue},
  pdfcreator={LaTeX via pandoc}}

\title{He--Kruczenski 有限原型：底层证明与适用范围}
\author{}
\date{2026-09-11}

\begin{document}
\maketitle

{
\setcounter{tocdepth}{2}
\tableofcontents
}
本册汇集已完成的底层推导、局部反例及严格数值证据的适用范围。它不宣称全部论文数值结果已经复现；当前物理运行和剩余项以同目录上一级的
PROGRESS.json、AUDIT\_WORKING\_ZH.md
及最终结果报告为准。原生1500盘、3123盘修补、3582盘修补及连续物理目标分别注明，不在模型之间搬用证明。

\newpage

\section{解析cardinal有限原型：独立推导归档}\label{ux89e3ux6790cardinalux6709ux9650ux539fux578bux72ecux7acbux63a8ux5bfcux5f52ux6863}

源：\path{results/runs/sequential_reproduction_20260910/DERIVATION_ZH.md}；SHA256=\nolinkurl{24180a41b18d1f58bc3ddf85d429d63df14b9fe652a652c0fad7c8153fb8e8bb}。此文保存该次完整推导，新增取样和端点修补的实现／见证另见\texttt{results/runs/physical\_consistency\_20260911}。不转移原生节点的可行性到新增网格。

\section{顺序主线：统一解析基及原参数联合问题}\label{ux987aux5e8fux4e3bux7ebfux7edfux4e00ux89e3ux6790ux57faux53caux539fux53c2ux6570ux8054ux5408ux95eeux9898}

\subsection{1.
先固定定义，不根据输出选择}\label{ux5148ux56faux5b9aux5b9aux4e49ux4e0dux6839ux636eux8f93ux51faux9009ux62e9}

本轮采用原文印出的两条手征残差不等式，各取四维 Euclidean
范数，epsilon=.002。保留打印 FESR 中心、四项 raw-absolute .002、原
uniform midpoint cutoff、FF squared caps 6e-5
和所有质量、耦合及凝聚输入。实际双密度 L4 的 \(B(M)=100[M^2+M(M+1)/2]\)
作为已声明的数值正则化保留，不能称它是原2309明示参数。

新模型取完整 M 模式 sine-cardinal 密度及其精确 Cauchy
变换。它是原文未指定插值的一种明确实现，不与旧 PV
混合近似相混。没有直接施加 T0=0、FF endpoint
等式、弹性饱和或峰质量条件。所有散射约束仍为原生节点的
\textbar S\textbar≤1。

\subsection{2.
一个函数产生所有行}\label{ux4e00ux4e2aux51fdux6570ux4ea7ux751fux6240ux6709ux884c}

令 phi\_j=pi(j+1/2)/M，D\_nj=(2-delta\_nM) sin(n phi\_j)/M。正弦正交给 D
S=I，因此转换保留全部 M 方向。对 q\_j(z)=sum\_n D\_nj
z\^{}n，b\_j=tan(phi\_j/2)/M，有 q\_j(-1)=-b\_j。未减除 Cauchy 核为
q\_j(z)+b\_j。

设 R 为任意双密度矩阵，Q 为对称双密度矩阵，Q\_ii=Cii、Q\_ij=Cij/2。令 b
为上述列向量。把未减除振幅写成 z 模式的转换为

\[
\begin{aligned}
c&=T_0+b^T(\sigma_1+2\sigma_2)+b^T(2R+Q)b,\\
a&=D(\sigma_1+2Rb),\\
d&=D(\sigma_2+(R^T+Q)b),\\
U&=DRD^T,\qquad V=DQD^T.
\end{aligned}
\]

于是整个振幅为

\[
A(s,t,u)=c+a^Tp(s)+d^T[p(t)+p(u)]
+p(s)^TU[p(t)+p(u)]+p(t)^TVp(u),
\quad p_n(s)=z(s)^n.
\]

同位旋由原交叉组合构造，部分波始终为 f=(1/4) integral P\_l
T。没有用不同的核分别替换物理和阈下取值。

\subsection{3.
已有解析数据的复用}\label{ux5df2ux6709ux89e3ux6790ux6570ux636eux7684ux590dux7528}

认证 registry 的原始行作用于 z\^{}n 模态系数；对称 tu 非对角基函数为
(z\_t\^{}n z\_u\textsuperscript{m+z\_t}m
z\_u\^{}n)/2。复用时实施上述转换的转置，得到完整 nodal C\_flat
行；不能把原行直接当 H，也不能丢掉 Nyquist 模式。

逐文件核对 registry 内 SHA256，按 canonical\_wave\_index
而非文件列表位置排布。所有旧代码仍保持非执行状态。新 provider
只读取数值区间数据。

独立检查采用旧 tip 全部3876系数，在预定 j=0,25,49 上比较 source
行积与当前 CardinalAmplitude 的严格独立角积分。九项全部相容，见
A1\_source\_audit/source\_audit.json。这确认坐标转换、1/4 投影、未乘
kappa 及 tu 对称因子；不转移旧 tip 的可行性。

\subsection{4.
自由列可以复用，但必须重新验证}\label{ux81eaux7531ux5217ux53efux4ee5ux590dux7528ux4f46ux5fc5ux987bux91cdux65b0ux9a8cux8bc1}

在原生节点，Im q\_j(s\_i)=delta\_ij，交叉道核为实。因此

\[
\operatorname{Im} f^0_0(s_i)=\tfrac32\sigma_{1,i}+\sigma_{2,i}+D^0_i(R,Q),
\qquad
\operatorname{Im} f^2_0(s_i)=\sigma_{2,i}+D^2_i(R,Q).
\]

其余分波的自由虚部列精确为零。代码先检查转换后的区间包含这些恒等值，再写成精确自由列；双谱列不这样处理。求解入口再验证这一布局，故
residual 消元的前提没有被移除。常数实部锚分别为5/2和1。

\subsection{5. 严格 IR
内点存在的构造}\label{ux4e25ux683c-ir-ux5185ux70b9ux5b58ux5728ux7684ux6784ux9020}

取
v\_j=sin(phi\_j)，T0=0，sigma1=sigma2=v/8，R=vv\textsuperscript{T，Q=vv}T/2。相应单核为
h(s)=1+z(s)，其色散密度 sin(phi)\textgreater0。物理交叉道 Legendre 投影

\[
J_l(s)=\tfrac14\int_{-1}^1 P_l(\mu)h(t)\,d\mu
\]

可写成正密度乘外区 Q\_l 的积分，故每个保留 l
都严格为正。对非零高波，虚部分别正比于 9 J\_l sin(phi)、3 J\_l
sin(phi)、J\_l sin(phi)（I=0、2、1）。两 S 波还具有正的单谱项。

因此有限所有盘可由足够小的共同正缩放满足： tau \textless{} min\_j 2 Im
f\_j/(kappa\_j \textbar f\_j\textbar²)。 同一缩放还能满足两 chi
球及实际双密度界。它证明存在一个新族的有限严格 IR
内点；浮点候选仍须重新检查，不能用径向缩放修复负虚部。

\subsection{6. 进入 UV 的次序}\label{ux8fdbux5165-uv-ux7684ux6b21ux5e8f}

先通过完整新 H 的 IR 内点检查，再 prepare-current 绑定同一 H。原生 FF 的
K+iI 不变，散射 Gram 使用新
H，所以电流的代数构造可复用，但必须重新保存模型身份。新联合见证通过之前，不做三点相移和分辨率扫描。

保存的 Float64 H
证书与实际解析行的区间核验分别报告。它们的差别是算子舍入误差，不再是两种不同的解析核。连续全能全自旋认证仍不作为原文有限原型的额外前置条件。

\subsection{7. 缺失分辨率的解析行：Gauss
余项而非经验阶数拟合}\label{ux7f3aux5931ux5206ux8fa8ux7387ux7684ux89e3ux6790ux884cgauss-ux4f59ux9879ux800cux975eux7ecfux9a8cux9636ux6570ux62dfux5408}

对未由认证数据覆盖的 M/L，直接投影同一有限多项式振幅。令
\(d=|s-4|\)，角变量换成 \(u=d(1-\mu)/2\)。物理区的交叉割线位于
\(u\le-4\) 和 \(u\ge d+4\)；阈下 \(0<s<4\) 时位于 \(u\le-s\) 和
\(u\ge d+s\)。因此分别取 gap=4 或 gap=s，不能混用。

把 {[}0,d{]} 分成从端点向中部几何增长的对称面板 {[}l,r{]}，中心 c、半宽
h。令 D=min(gap+l,gap+d−r)，在面板标准坐标中选 Bernstein 椭圆的实半轴 a
满足 h(a−1)≤D/2。代码取向下舍入的 a≤1+D/(2h)，再取向下舍入的
\(\rho=a+\sqrt{a^2-1}>1\)。整个椭圆留在两个割线之间，故交叉 conformal
map 均满足
\textbar z\textbar≤1。于是任意保留幂次及其两个交叉道乘积的模均≤1，这个界不依赖
M。

若面板上 \(|\mu|\le R\)，Legendre 积分表示给出 \[
|P_\ell(\mu)|\le\bigl(R+\sqrt{R^2+1}\bigr)^\ell=:B_\ell.
\] 对椭圆内解析且模≤B的函数，Chebyshev 系数满足
\(|a_k|\le2B\rho^{-k}\)。用次数 2N−1 的截断多项式，余项一致界为
\(2B\rho^{-2N}/(1-\rho^{-1})\)。N点 Gauss
公式对该多项式精确，积分测度和正 Gauss 权重的总质量均为2，故 \[
\left|\int_{-1}^1g-Q_Ng\right|
\le {8B\rho^{-2N}\over1-\rho^{-1}}.
\] 本论文分波还含 1/4，角换元给 dμ=2du/d，面板换元给 du=h
dx；所以代码每面板加的是 \[
{4h\over d}\,{B_\ell\rho^{-2N}\over1-\rho^{-1}},
\] 而不是直接照搬8B。这给每个单模式积分 J 和双模式积分 U
的显式误差球。之后的同位旋组合、DST转置和未减除变换全部用区间矩阵运算传播误差，最后
Float64 H 的 radius\_upper
还包含存储舍入误差。N≥64是在看新曲线之前固定的计算选择，不是对曲线调参。

对称积分恒等式 Uji=(−1)\^{}ell Uij
在同一换元下成立；奇数ell时对角为零。只有原包络包含0时才写成精确零。原生自由虚部列也只在包络包含已证明的恒等值后精确化。独立角积分测试覆盖
S、P 以及非主分波的同一完整随机系数振幅。

一般精确/区间能量不按 float 相等关系吸附到原生节点；只有 CLI
明确约定的浮点节点别名或显式原生节点对象才使用缓存。非有限能量和跨越阈值的区间拒绝。这些接口条件保证``原生节点缓存''和``一般能量包络''不会互相冒充。

\subsection{8.
中心收敛、支持误差与物理曲线是三个问题}\label{ux4e2dux5fc3ux6536ux655bux652fux6301ux8befux5deeux4e0eux7269ux7406ux66f2ux7ebfux662fux4e09ux4e2aux95eeux9898}

对固定可行域，写障碍目标 \[
\Psi_\mu(z)=\Phi(z)-c^Tz/\mu,\qquad
\Phi=-\sum_j\log\Delta_j-w_\chi\sum_a\log(\epsilon^2-\|r_a\|_2^2)+\Phi_\rho(\rho/B)+\Phi_{\rm current}.
\] 这里实现的是保留原L4可行域的辅助变量消元障碍，不能误写成简单的
−log(1−sum r\_i\^{}4)。令 r\_i 为实际密度除以 B， \[
\Phi_\rho(r)=\min_{v_i>r_i^2,\;v^Tv<1}
\{-\sum_i\log(v_i-r_i^2)-\log(1-v^Tv)\}.
\] 存在这样的 v 当且仅当 sum r\_i\^{}4\textless1。记
a\_i=r\_i²、d\_i=v\_i−a\_i、δ=1−sum v\_i²，内层驻点满足 δ=2v\_i
d\_i，因此 \[
v_i={a_i+\sqrt{a_i^2+2\delta}\over2},\qquad
F(\delta)=\delta+\sum_i v_i(\delta)^2-1=0.
\] F(0)=sum a\_i²−1\textless0，F(1)\textgreater0，且 F′=1+sum
v\_i/sqrt(a\_i²+2δ)\textgreater0，故唯一正根存在。代码的标量Newton步正是对这个F求根，不是另设密度正则器。消元后的梯度为
g\_i=2r\_i/d\_i=4r\_i v\_i/δ。再隐式微分可得 \[
\nabla^2\Phi_\rho=\operatorname{diag}
\left({2\over d_i}+{8a_i\over\delta+2d_i^2}\right)+uu^T,
\quad
u_i={-4r_i v_i\over(\delta+2d_i^2)\sqrt{1+\sum_j\delta/(\delta+2d_j^2)}}.
\]
末项确实是正的秩一项。对完整辅助Hessian作Schur消元不能据直觉把最后的显式表示改成负号。恢复未归一化密度坐标时再乘1/B或1/B²；代码的特征矩阵与独立有限差分检查使用相同规则。

每个权重均为正。把 wχ 从1变为750不改变任何可行不等式，也不改变投影目标
c；它改变有限μ时选择的中心路径。选取750的规则是
n\_disks/2，与已使用的联合求解器一致，不由相移决定。梯度、Hessian、沿射线导数和对偶权重必须同时乘相应权重，不能只改停止量。

固定 \(Xz=x_{\rm ref}\) 时，以 \(V=(-X_{1:}/X_0;I)\) 参数化切空间（这里
X0≠0）。Newton 方程为 \[
(V^T\nabla^2\Psi V)d=-V^T\nabla\Psi,
\qquad \Delta z=Vd,
\] 因此全部密度方向仍在该仿射截面中优化。法向乘子由完整梯度沿 X
的分量恢复，形成全局支持方向
\((\lambda,n_y)\)。径向中心换元不能在这条固定截面迭代中直接缩放 z。

这里的联合中心是消去自由高能电流谱后的约化问题中心，仍保留全部3876个振幅方向。原始4076变量中的14个高能谱对角元既无矩上界也不进目标；若对它们直接使用完整−logdet障碍，沿R→∞会下降，不能宣称原始全变量问题有唯一有限障碍中心。严格散射内点上用Schur公式消去这些自由谱，再将约化中心拓展为满足全部原式的完整点，才是当前算法与证书的准确关系。

固定μ\textgreater0时，这个约化中心在严格内点域中是唯一的。上面的密度Hessian具有严格正对角加半正定秩一项，因此任何非零密度方向都有正曲率。若剩余自由方向使散射障碍的Hessian二次型为零，则所有原生分波的实、虚部变化都为零；先由S2的自由虚部列得δσ2=0，再由S0得δσ1=0，最后由S0实常数列5/2得δT0=0。因此IR
Hessian正定，固定xref的切空间限制也仍正定。

对联合约化问题，先由上述论证得δC=0。低能正定Gram的−logdet曲率为 tr(G⁻¹δG
G⁻¹δG)，为零当且仅当δG=0，因而所有低能ImF及受矩约束的R变化为零。高能FF盘的正二次曲率又要求各高能ImF变化为零。故联合约化Hessian也正定。密度/散射界使C有界，正FESR权重和FF上界使保留电流变量有界；在严格内点存在的截面上，障碍在相应边界发散，所以每个固定μ的最小值存在且唯一。这里不能只从无截面的Phase
I推断任意固定x截面有内点。

本轮ref/mid所需的存在前提可由已有两点具体满足：D3\_analytic\_primal\_lifted与E1\_tip\_verified的约化原约束全部严格，其解析x分别约.040853448和.099039809，跨越xref及xmid。对任一中间x取精确权重
t=(x−x\_D3)/(x\_tip−x\_D3)，完整点的精确凸组合即在目标截面严格可行。两点86个低能Gram均已由主子式严格判正；部分高能Gram只有PSD证明，但这些自由高能R不参与约化障碍，必要时各加任意正量即可严格提升。详见\href{../../sequential_reproduction_20260910/PROOF_SCOPE_REVIEW_ZH.md}{独立适用性审阅}。因此本轮这两个截面的中心存在唯一性前提已补齐，而不是假设所有截面自动满足。此结论不证明μ→0的极值幅度唯一，也不证明ρ极点或物理曲线误差。

收敛中心要求 Newton
decrement、实际线性残差和残余梯度的支持宽度共同达标。早先``一次Newton后目标gap已小''只说明那一候选的投影间隙，不能证明它已到中心。本轮
\texttt{-\/-require-center}
保留真正收敛的完整中心作为代表，停止不再只由旧对偶和新非中心primal之间的gap触发。

已完成 tip 的6次精修给
decrement≈1.82e−14、支持gap≈7.44e−5。ref/mid也分别在其固定截面上完成中心与支持。最后两层μ间，ref/mid
的 P1 最大模π相位变化约0.256°/0.611°；这只是实测变化，不能当成 μ→0
的误差定理。支持gap只控制线性目标，不控制所有分波或保证极值振幅唯一。末端极小的浮点可行恢复缩放、以及自由高能电流谱的Schur提升均另行记录，不能隐称保存点与实数中心逐位相等。

IR
的权重1路径曾在极薄手征余量处停滞；直接重算状态发现最大分波坐标漂移约2.3e−14，手征余量漂移约3.6e−20，不支持``缓存状态明显漂移''的解释。采用相同750权重规则后，绿色IR中心收敛，支持gap约5.44e−4，对应纵向gap约1.09e−6；它又通过全部1500解析散射行、两条χ和L4的区间核验。

\subsection{9. 本轮已经闭合的 IR→UV→ρ
证据链}\label{ux672cux8f6eux5df2ux7ecfux95edux5408ux7684-iruvux3c1-ux8bc1ux636eux94fe}

在固定 xref
的绿色IR上侧中心，P1相移仍很小，未出现90°上穿或具有ρ强度的共振峰。完整C先保存于
C2\_green\_selection，之后才求值。三个UV代表同样在选点之前固定目标和规则，全部4076原变量的原约束及解析primal核验均通过。其90°读数为
tip约803.49、mid约701.26、ref约696.63MeV；原生散射强度主峰分别为792.14、680.41、680.41MeV。没有加入ρ质量、宽度或相位拟合条件。

固定绿色IR的整个C，只求电流扩展，原保存有限算子的Farkas上界约−9.10961718\textless0。证明形式是：所有电流Gram
PSD、FF锥和四个矩区间的非负组合，加上已由这些原约束推导的变量残差上界，推出零目标≤负数。PSD对偶修正经Arb主子式检查，不能只依赖求解器的
PrimalInfeasible
字样。此证书排除的是这个完整C，而不是其二维投影、整个联合可行域或所有``没有ρ''的函数。

Gram Schur 恒等式（\(\Delta=1-|S|^2>0\)）为 \[
R\ge |\mathcal F|^2+{|\mathcal F-S\mathcal F^*|^2\over\Delta}.
\]
F(0)=1与高能FF界限制形状因子的解析变化，UV矩限制可用电流谱预算，偏离散射/FF相位一致性的项消耗额外谱预算。因此原UV信息可以排除IR代表并推动出现共同中能结构。上述严格不等式没有自动要求Gram秩1，也没有证明唯一ρ极点或正确质量；矩约束本身也不是整个谱的唯一性条件。

所以本轮已复现``加入UV信息后出现ρ式结构''的有限数值机制；三代表中两条共振偏低且非弹性较强的
major
差异仍在。中心精修和共同解析原型修复没有消除这一差异，不能将其重新归类为小的Newton未收敛误差。新M/L、完整区域及收缩非对称仍须按顺序检查，不能用旧PV状态代替。

固定C证书现已用768-bit重新计入全部相关解析源行的误差球，并改用精确原生K、运动学、打印十进制矩、raw容差和hard权重。新严格上界为
−9.109617136909641，仍明确为负；相对原有限H上界仅变化约4.12e−8。证据见
CD\_fixed\_current\_analytic/source\_audit.json。这里的排除已不依赖忽略H的存储/积分误差。

\subsection{10.
新增的节点间反例}\label{ux65b0ux589eux7684ux8282ux70b9ux95f4ux53cdux4f8b}

本轮后续直接求值和18项独立积分发现三UV振幅在ρ附近有严格η\textgreater1，详见
\href{../../sequential_reproduction_20260910/HELD_INTER_NODE_ISSUE_ZH.md}{独立汇总}。第9节的机制结论严格限于原生有限程序；它不能被解释为已经得到连续物理上合格的ρ相移。按照用户对这类反例的要求，网格加密暂列独立修补方案，未并入原基线。

\subsection{11.
两条向量矩如何容许偏低的ρ}\label{ux4e24ux6761ux5411ux91cfux77e9ux5982ux4f55ux5bb9ux8bb8ux504fux4f4eux7684ux3c1}

原文Eq.(3.72--74)写原始离散积分与QCD值的绝对差；Eq.(2.56)打印的是除以s0\^{}(n+2)后的数值，必须先乘回。当前逐项绝对容差.002有印式依据；原作者是否另做范数打包或代码缩放仍未公开，不能因为别的口径更接近相位就改用它。原文见
\texttt{references/2309.12402v3-source/prd\_submission\_2.tex:1022–1040}。

令正的离散测度权重为 a\_i=(π/M)s′\_i，保留s\_i≤s0。向量谱R\_i≥0，因此 \[
I_{-1}=\sum_i a_iR_i/s_i>0,\qquad
\mu_i={a_iR_i/s_i\over I_{-1}},\quad\sum_i\mu_i=1,
\qquad \overline{E^2}=m_\pi^2\sum_i\mu_i s_i=m_\pi^2{I_0\over I_{-1}}.
\] 这里是以R/s为权重的能量平方均值，不是任意峰位或极点质量的公式。

打印中心值转回原始矩后是
I−1=0.004228045714\ldots、I0=0.145711206997\ldots。同样的绝对.002误差，对前者相当于47.3\%，对后者约1.37\%。只由这两条区间推出的加权均方能标范围约为0.6725--1.1399GeV，中心值为0.8219GeV。这既给出为何UV引入一个中能尺度，也说明它单独并不固定770MeV，更不保证谱有窄峰。

当前三个UV点实际使用的向量矩如下：

{\def\LTcaptype{none} % do not increment counter
\begin{longtable}[]{@{}lrrrr@{}}
\toprule\noalign{}
点 & I−1 & I0 & 加权均方能标 MeV & 直接加点90°读数 MeV \\
\midrule\noalign{}
\endhead
\bottomrule\noalign{}
\endlastfoot
tip & .00443308288 & .14771073369 & 808.13 & 802.01 \\
mid & .00590775693 & .14771083475 & 700.04 & 697.13 \\
ref & .00609361320 & .14771082229 & 689.28 & 688.38 \\
\end{longtable}
}

三点的I0均靠近允许上界，I−1的变化显著改变了这个谱能标。数据来自
\texttt{CE\_phases/phases.json}；没有用上述比值反向设定任何振幅或峰。它解释了为何当前低峰仍通过给定的UV矩，而不是把它归为Newton误差。

Gram条件进一步限制相位：R≥\textbar ℱ\textbar²+\textbar ℱ−Sℱ*\textbar²/(1−\textbar S\textbar²)。若IR的S几乎为1，而由归一化、解析性和高能界产生的F有较大虚部，则相位失配会消耗很大的有限谱预算；固定IR振幅的负Farkas界是这一冲突的全变量证据。改变S的相位或增加非弹性都可降低这项代价，原不等式没有强制二π饱和。

在δ=90°且F相位一致时，R的Schur下界为2\textbar ℱ\textbar²/(1+η)，对应二π谱比例最多达到(1+η)/2；散射强度则为(1+η)²/4。因此一个很高的FF/电流谱峰可以伴随较弱、非弹性很强的散射峰。当前原生三点的η差异正需要保留为major物理差异，不能用``Gram近饱和''代替弹性或完整ρ复现。

\subsection{12.
谱矩预算的必要作用：一个直接构造}\label{ux8c31ux77e9ux9884ux7b97ux7684ux5fc5ux8981ux4f5cux7528ux4e00ux4e2aux76f4ux63a5ux6784ux9020}

可以严格说明为什么只有高能FF上界还不够。此处仅作数学上的约束删减比较，不改生产问题：删去四条FESR，保留任意一个原生散射盘严格可行的IR振幅、F(0)=1、原FF上界和电流Gram。对M≥4，当前完整FF多项式空间包含
\[
F_1(z)=(1+z)^2,\qquad F_0(z)=(1+z)^4.
\]
因为z(0)=0，两者都满足归一化；实系数确保实解析性，且这些系数可以由满秩的原生正弦变换完整表示。物理割线上z=e\^{}\{iφ\}、s=4/cos²(φ/2)，故\textbar1+z\textbar²=16/s。用原运动学因子、β=√(1−4/s)，有
\[
k_1^2|F_1|^2={2\beta^3\over3\pi^5s}
\le{2\over3\pi^5s_0}<3\times10^{-5},
\] \[
k_0^2|F_0|^2={768\beta\over\pi^5s^4}
\le{768\over\pi^5s_0^4}<2m_q^2(6\times10^{-5}).
\]
这里s0=(1.2/.14)²、mq=(.004+.0073)/(2×.14)。两个保守上界分别约2.96519×10⁻⁵和8.61363×10⁻⁸，严格低于原向量3×10⁻⁵及标量1.95444×10⁻⁷上界，并且对所有s≥s0成立。没有使用ρ质量或相位。

有限个原生节点上Δ\textgreater0，F有限，故每个Schur下界都有限。删去矩预算后，各R可独立取在该下界之上，得到全部电流Gram正定。因此这个删减问题可扩展\textbf{任意严格IR振幅}；高能FF界本身不能排除该IR代表。加入FESR后，各低能R不能再任意增大，第9节负Farkas界证明固定绿色IR恰在这个有限预算下被排除。这给出谱矩在机制中不可替代的直接论证；它仍不证明矩预算必然产生唯一或正确的ρ。

\subsection{13.
当前解析族的有限紧性与支持：独立重证}\label{ux5f53ux524dux89e3ux6790ux65cfux7684ux6709ux9650ux7d27ux6027ux4e0eux652fux6301ux72ecux7acbux91cdux8bc1}

令r为实际双密度独立条目（存储的对称非对角C必须除2），设\textbar\textbar r\textbar\textbar₄≤B。第4节原生恒等式已经对新解析行重新验证；写D0i、D2i为对应的实际密度行，则由散射盘有0≤uIi≤2/κi，进而
\[
|\sigma_{2i}|\le {2\over\kappa_i}+B\|D_{2i}\|_{4/3},\qquad
|\sigma_{1i}|\le {4\over3\kappa_i}+{2B\over3}\|D_{0i}-D_{2i}\|_{4/3}.
\] 任取一个S0原生实部行Re f00=5T0/2+a·σ+Dr，用\textbar Re
f\textbar≤1/κ，得到 \[
|T_0|\le {2\over5}\left({1\over\kappa}+\sum_j|a_j|\overline\sigma_j+B\|D\|_{4/3}\right).
\]
固定M时所有κi\textgreater0、行元素有限，因此全部C有界；闭约束与零振幅给出非空紧集，所有IR线性支持均取得最大值。这是以新族的原生恒等式为前提的证明，不转移PV或有限sine的秩结论；L4单独不能约束未直接正则化的σ及T0。H有3011行、3876列，本来也不需要、也不可能满列秩。

对κ已吸收入的坐标(x,y)=(κRe f,κIm f)，盘中心为(0,1)、半径1，支持函数为
\[
\sup_{x^2+(y-1)^2\le1}(a x+b y)=b+\sqrt{a^2+b^2}.
\] 若系数直接乘未缩放的Re f、Im
f，整个右侧须除κ。密度球的支持是B\textbar\textbar r\_dual\textbar\textbar₄/₃；所有自由列须精确消去或有独立有界残差控制。当前验回检查原生自由列身份后再消元，未要求整个H可逆。保存Float64
H的对偶支持与实际解析H±E的primal仍分开声明；固定C的电流排除另有计入解析误差的证书，不能据此宣称全部联合对偶已升级为解析算子证书。

联合问题中低能R由正矩权重界定、所有ImF由低能Gram或高能FF界定，但高能自由R无上界，完整4076维集并不紧。消去高能R所得紧松弛集可能含Δ=0且不满足Watson范围条件的边界点；与一个严格点作凸混合可令Δ\textgreater0，再有限Schur提升。故其支持值等于原投影闭包的上确界，不自动保证完整联合极值在有限高能R处取得。当前每个交付点仍独立补回并通过全部原式，未把形式闭包点当物理解。

\newpage

\section{有限解析族的领先高能幺正条件}\label{ux6709ux9650ux89e3ux6790ux65cfux7684ux9886ux5148ux9ad8ux80fdux5e7aux6b63ux6761ux4ef6}

范围：当前M模式未减除cardinal振幅，T0=0；所有系数实，R任意、Q对称，Q的非对角原存储C等于2Q。结论对每个固定分波成立，严格条件可给任意有限保留分波集合的最终高能幺正性；不提供无限自旋统一界，也未给出有限
crossover 能量。

\subsection{极限与积分换序}\label{ux6781ux9650ux4e0eux79efux5206ux6362ux5e8f}

写h\_j=q\_j+b\_j，q\_j(−1)=−b\_j，令a\_j=4q'\_j(−1)。物理割线上

\[
1+z(s)=(8+4i\sqrt{s-4})/s,
\quad\sqrt{s}\,\operatorname{Im}h_j(s)\to a_j,
\quad\sqrt{s}\,\operatorname{Re}h_j(s)\to0.
\]

对t=−(s−4)(1−x)/2，有

\[
\sqrt{s}\,h_j(t)\to {\sqrt2 a_j\over\sqrt{1-x}}.
\]

因为\(h_j(z)=(1+z)\widetilde q_j(z)\)，有限多项式在闭单位盘上一致有界。s≥8时有√s\textbar1+z(t)\textbar≤8/√(1−x)，故右端存在与s无关的可积控制函数。两个交叉核之积乘s也被常数/√(1−x²)控制，可积分。因此相应积分换序合法，而非逐角极限的形式代入。

Rodrigues公式或Beta积分给

\[
\int_{-1}^{1}{P_l(x)\over\sqrt{1-x}}dx={2\sqrt2\over2l+1}.
\]

具体地，对Rodrigues公式分部积分l次，每一边界项在x→1时均为O(√(1−x))、在x→−1时至少有一个(1+x)因子，因此为零。剩余积分为(1/2)\_l/(2\^{}l
l!)乘∫(1−x)\textsuperscript{(-1/2)(1+x)}l
dx。换元y=(1+x)/2后，用B(l+1,1/2)，得到√2 Γ(l+1/2)/Γ(l+3/2)=2√2/(2l+1)。

用J\_lj=(1/4)∫P\_l h\_j(t)dx，得到√s
J\_lj→a\_j/(2l+1)。两个交叉核乘积的积分为O(1/s)，故不贡献分波实部的1/√s领先项；它们在物理区为实，也不贡献虚部。

\subsection{由交叉组合得到五项条件}\label{ux7531ux4ea4ux53c9ux7ec4ux5408ux5f97ux5230ux4e94ux9879ux6761ux4ef6}

定义

\[
\alpha_1=a\cdot\sigma_1,\quad\alpha_2=a\cdot\sigma_2,
\qquad r=a^TRa,\quad q=a^TQa.
\]

两条S波虚部有

\[
\sqrt{s}\,\operatorname{Im}f_0^0\to\tfrac32\alpha_1+\alpha_2,
\qquad\sqrt{s}\,\operatorname{Im}f_0^2\to\alpha_2.
\]

对非零保留分波，令d\_l=2l+1，则

\[
\sqrt{s}\,\operatorname{Re}f_l^I\to{2A_I\over d_l},
\qquad s\,\operatorname{Im}f_l^I\to{2K_I\over d_l},
\]

\[
(A_0,A_1,A_2)=(\alpha_1+4\alpha_2,\alpha_1-\alpha_2,\alpha_1+\alpha_2),
\qquad(K_0,K_1,K_2)=(4r+q,r-q,r+q).
\]

例如I=0的直接×交叉R项给虚部8r/d\_l，Q项给2q/d\_l；I=1给2(r−q)/d\_l。因κ→π，Δ=2κImf−κ²\textbar f\textbar²满足

\[
s\Delta_l^I\to {4\pi\over d_l^2}(d_lK_I-\pi A_I^2),\qquad l>0.
\]

因此全高能幺正性的必要条件是

\[
3\alpha_1+2\alpha_2\ge0,\quad\alpha_2\ge0,
\] \[
5(4r+q)\ge\pi(\alpha_1+4\alpha_2)^2,
\quad3(r-q)\ge\pi(\alpha_1-\alpha_2)^2,
\quad5(r+q)\ge\pi(\alpha_1+\alpha_2)^2.
\]

最低非零偶波l=2和奇波l=1最严格；满足它们即满足其它固定更高波的领先必要条件。若五式全部严格，由极限定义存在足够大能量，使任意给定有限保留波集合全部Δ\textgreater0。等号时次领先项仍可能改变结论，不能用≥0冒称严格的最终高能证书。

\subsection{凸性、种子与对偶}\label{ux51f8ux6027ux79cdux5b50ux4e0eux5bf9ux5076}

r和q对双密度系数线性，α对单密度线性。后三式都是线性项不小于线性项平方，构成凸集合；不是按ρ峰加入的经验约束。可用五个形式2y−x²≥0表示：前两个取x=0，后面取x=√(π/d\_l)A\_I、y=K\_I/2。正的尺度变换x→x/√w、y→y/w不改集合。

对闭抛物域2y≥x²，支持函数为

\[
\sup (u x+v y)=-{u^2\over2v}\quad(v<0),
\]

v=0时仅u=0给有限值0，其余无有限支持。由−log(2y−x²)的梯度产生v=−2μ/(2y−x²)\textless0。验回应把实际u、v乘原解析行计入站立性残差，并加上述支持常数；不能沿用圆盘的支持公式。

正的IR种子已存在：取σ1=σ2=λ sinφ/8、R=λ vvᵀ、Q=λ
vvᵀ/2、T0=0。由于a·v=4，α1=α2=λ/2、r=16λ、q=8λ，对足够小λ\textgreater0五式均严格。再结合有限所有散射盘及χ/L4的小缩放构造，可作为完整Phase
I的硬IR内点。它不证明加入UV后必定可行；UV仍须完整求解、原式及解析核验。

端点切片的第一个完整联合见证确实违反后三式，数值包络见\path{results/runs/physical_consistency_20260911/endpoint_witness_analytic/source_audit.json}。这是当前点的严格高能障碍，不是整个端点切片无解。新增五式后必须另立模型记录并重新求见证。

\subsection{独立有限能量检查}\label{ux72ecux7acbux6709ux9650ux80fdux91cfux68c0ux67e5}

\path{results/runs/physical_consistency_20260911/TAIL_LIMIT_CONTROL.json}使用手工定义的M3首谐波v=(1/2,1,1/2)、λ=1/1024，完整解析投影分别在s=10\textsuperscript{6、10}8求值。α1=α2=λ/2、r=16λ、q=8λ由h=1+z直接给定，未从同一导数矩阵计算期望值。六个保留波的相应缩放虚部接近所推导极限，误差随能量提高明显下降。这只是独立数值交叉检查；上述控制收敛与积分计算才是极限证明，不由两个能量样本代替。

\newpage

\section{有限解析族的阈值幺正必要条件}\label{ux6709ux9650ux89e3ux6790ux65cfux7684ux9608ux503cux5e7aux6b63ux5fc5ux8981ux6761ux4ef6}

本结论针对本仓库的未减除有限 cardinal
解析族；不是原文另行给定的输入。目前只推导并独立检查，\textbf{没有加入生产优化约束}。高能必要条件成立不能替代这里的阈值条件，也不能替代中间能区的幺正性。

设 d=s−4↓0，h\_j(z)=b\_j+∑\emph{\{n=1\}\^{}M
D}\{nj\}z\^{}n，e\_j=h'\_j(1)。物理割线上 z=(8−s+4i√(s−4))/s，所以

\[
\operatorname{Im}h_j(s)=e_j\sqrt d+O(d^{3/2}).
\]

实系数、有限多项式保证余项在任意固定有限核集合上一致。对交叉变量
t=−d(1−x)/2，h\_j(t)在t=0附近解析，其Taylor级数在全部x∈{[}−1,1{]}上一致收敛。因此逐项投影合法。

由 t=16z/(1+z)\^{}2 和 Lagrange 反演，对 l≥1有

\[
c_{lj}:=[t^l]h_j(t)=16^{-l}\sum_{n=1}^{\min(M,l)}D_{nj}\frac nl {2l\choose l-n}.
\]

定义
J\_\{lj\}=(1/4)∫\_\{−1\}\^{}1P\_l(x)h\_j(t)dx。Legendre正交性消掉所有次数小于l的项，而

\[
\int_{-1}^1P_l(x)(1-x)^l dx
=(-1)^l\frac{2^{l+1}(l!)^2}{(2l+1)!}.
\]

此恒等式直接由 Rodrigues
公式分部积分l次得到；(x²−1)\^{}l使所有边界项为零。因此

\[
J_{lj}=g_{lj}d^l+O(d^{l+1}),\qquad
g_{0j}=b_j/2,\qquad
g_{lj}=\frac{(l!)^2}{2(2l+1)!}c_{lj}\quad(l\ge1).
\]

以下仅讨论允许的分波：I=0、2取偶数l，I=1取奇数l；禁戒奇偶的振幅恒等为零，不能套用下面的非零波系数。所有交叉核在物理区为实。按原交叉组合，令R为非对称双密度矩阵、Q为对称双密度矩阵（存储的非对角系数是2Q），得到

\[
\begin{aligned}
B_l^0&=6e^TRg_l+2g_l^TRe+2e^TQg_l,\\
B_l^1&=2g_l^TRe-2e^TQg_l,\\
B_l^2&=2g_l^TRe+2e^TQg_l.
\end{aligned}
\]

对l=0，第一式还应加(3/2)e·σ₁+e·σ₂，第三式还应加e·σ₂。于是 Im
f\_l\textsuperscript{I=B\_l}I
d\textsuperscript{\{l+1/2\}+O(d}\{l+3/2\})。对l\textgreater0，正交性同时给
Re
f\_l\textsuperscript{I=O(d}l)；对S波令A\_I=f\_0\^{}I(4)，它是完整C的实线性泛函。

由于κ=π√d/2+O(d\^{}\{3/2\})，原幺正裕量Δ=2κIm
f−κ²\textbar f\textbar²满足

\[
\lim_{d\downarrow0}\frac{\Delta_0^I}{d}
=\pi B_0^I-\frac{\pi^2}{4}A_I^2\quad(I=0,2),
\qquad
\lim_{d\downarrow0}\frac{\Delta_l^I}{d^{l+1}}
=\pi B_l^I\quad(l>0).
\]

因此阈值邻域幺正必需
B₀ᴵ≥πAᵢ²/4，以及每个非零保留波Bₗᴵ≥0。S波条件是凸抛物域，高波条件线性。对任意固定有限波集合，若这些条件全部严格，则由上述极限存在共同的δ\textgreater0，使4\textless s\textless4+δ上的全部保留波幺正。等号情形必须考察更高阶；本结论不给出δ数值，也不处理无限自旋一致性。

\subsection{独立首谐波检查及适用范围}\label{ux72ecux7acbux9996ux8c10ux6ce2ux68c0ux67e5ux53caux9002ux7528ux8303ux56f4}

取M=3、v=(1/2,1,1/2)、λ=1/1024、σ₁=σ₂=λv/8、R=λvvᵀ、Q=λvvᵀ/2、T₀=0。此时∑v\_jh\_j=1+z，不需要上述D矩阵就可直接展开：

\[
g_l=\frac1{2(2l+1)(l+1)16^l},\quad
A_0=12\lambda,\quad A_2=\tfrac92\lambda,
\quad B_0^0=77\lambda/16,\quad B_0^2=13\lambda/8,
\]

非零波有B\_l\textsuperscript{0=9λg\_l、B\_l}1=λg\_l、B\_l\^{}2=3λg\_l。
单一CLI计算
\path{results/runs/physical_consistency_20260911/threshold_limit_control_evaluate}
在 s=4.000001、4.00000001
直接投影六个波；实际binary64能量作为精确输入，并保存解析包络。它用于检查极限系数和幂次，不是QCD联合可行见证。对应数值比较见
\path{results/runs/physical_consistency_20260911/THRESHOLD_LIMIT_CONTROL.json}：两能量上的最大相对虚部系数误差约6.0×10⁻⁷、6.0×10⁻⁹；裕量系数误差约7.25×10⁻⁷、7.25×10⁻⁹，与O(d)余项一致。

这里的证明解释为何极小的高波负裕量也可能是严格物理违例：其自然量级本来就是d\^{}(l+1)。不能将统一的绝对浮点容差当作物理可行性标准。

\newpage

\section{同一解析模型的IR截面、可行中心与原式支持证明}\label{ux540cux4e00ux89e3ux6790ux6a21ux578bux7684irux622aux9762ux53efux884cux4e2dux5fc3ux4e0eux539fux5f0fux652fux6301ux8bc1ux660e}

本节针对新增的 \texttt{ir.py}
通路。它与加强后的UV模型使用完全相同的散射振幅、3123个H行、T0=0、五个领先高能条件、手征球与实际双密度L4；IR不含FF、电流谱和FESR。M3/L2控制仅验证实现，不能替代M50/L10的物理复现。

\subsection{1.
可逆坐标与有限非空域}\label{ux53efux9006ux5750ux6807ux4e0eux6709ux9650ux975eux7a7aux57df}

设完整原系数为C，p=1+2M+M²+M(M+1)/2。以D表示对rho2非对角存储项乘2、其他项不变的对角矩阵，则u=D⁻¹C的双密度段是实际密度。已有absorptive变换仅将前2M组单谱换成原生Im
f00及Im f20，并保留T0及每一个双密度方向；其逆记作T。

原生盘

\[
(\kappa\Re f)^2+(1-\kappa\Im f)^2\le1
\]

严格推出0≤κIm f≤2。在T0=0仿射切片中，数值变量取

\[
\xi=(\kappa_i\Im f^0_0(s_i),\kappa_i\Im f^2_0(s_i),r),
\qquad C=D\,T\,\operatorname{diag}(1,\kappa^{-1},1)(0,\xi)^T.
\]

这是p−1维切片上的线性双射；密度维数不减，T0由代数恒等式精确等于零。前2M变量的{[}0,2{]}界来自原盘，不是新增物理假设。此解释不能移用到仍保留自由T0的数值变量。

变量属于有界闭集{[}0,2{]}\^{}(2M)×\{‖r‖₄≤B\}；再与所有盘、手征球、五个凸高能条件及固定x仿射截面相交仍为紧集。因此非空时任何线性目标取得最大值。这里仅为IR有限截面；联合UV中未受矩约束的高能电流谱仍须另论。

非空性按算子身份分别论证：闭式第一谐波证明精确解析IR函数族非空；它的严格性不能未经误差检查直接转移到保存的float64
H。对于本次实际有限H及具体物理xref，若已有在同一H上独立核验的UV严格点满足x\textgreater xref\textgreater0，则α=xref/x∈(0,1)的完整C缩放为该IR截面的严格初值。所有IR不等式均保留：盘和高能抛物裕量满足

\[
2\alpha y-\alpha^2x^2
=\alpha(2y-x^2)+\alpha(1-\alpha)x^2\ge0,
\]

L4及手征范数随α缩小。只缩放散射C；归一化FF不满足同质缩放，所以这一步绝不声称得到一个新的UV点。

\subsection{2.
五个领先高能条件的障碍与对偶}\label{ux4e94ux4e2aux9886ux5148ux9ad8ux80fdux6761ux4ef6ux7684ux969cux788dux4e0eux5bf9ux5076}

由独立高能推导得到的原系数行记作R\_a、I\_a，a=1,\ldots,5，约束为

\[
2I_aC-(R_aC)^2\ge0.
\]

其中两个R行恒零，给两个线性条件；其余三个是凸抛物约束。换元和任意固定正数w\_a仅改成

\[
x_a=R_aC/\sqrt{w_a},\quad y_a=I_aC/w_a,
\qquad 2y_a-x_a^2\ge0.
\]

域没有改变。障碍−log(2y−x²)的方向导数为(2x dx−2dy)/(2y−x²)，二阶型分解为

\[
\left(\frac{2x\,dx-2dy}{2y-x^2}\right)^2
+\frac{2(dx)^2}{2y-x^2}.
\]

这直接给实现的两类Hessian特征行，不从数值曲线反推。

对u x+v y的支持函数，v\textless0时

\[
\sup_{2y\ge x^2}(ux+vy)=-\frac{u^2}{2v};
\]

v=0仅u=0有限，否则支持无界。障碍协向量u=2μx/(2y−x²)、v=−2μ/(2y−x²)满足有效符号。回到物理原行的协向量为u/√w和v/w；原C残差中减去这两项，支持常数不变。最终
\texttt{audit\_ir}
直接用原R、I和Arb重新计算残差、支持常数及五个primal裕量。

\subsection{3.
Newton小误差不是支持证明}\label{newtonux5c0fux8befux5deeux4e0dux662fux652fux6301ux8bc1ux660e}

设数值站立性残差为g，按自由/密度段分成g\_f、g\_r。任意两个IR可行点之差满足

\[
|g\cdot(\xi_1-\xi_2)|
\le2\|g_f\|_1+2B\|g_r\|_{4/3}.
\]

这给出有全域意义的梯度宽度。Newton
decrement只度量局部Hessian几何；大B和很小的某些曲率可以使decrement小而上述宽度仍大。新的中心停止同时检查decrement、线性方程残差及梯度宽度。即使这些数值指标通过，最后仍需回到原H的独立对偶支持和primal核验；任何一个数值指标都不替代它们。

实际密度障碍不是单独的−log(1−Σρ\_i⁴)。令ρ=r/B；实现使用提升域

\[
 v_i>\rho_i^2,\quad \sum_i v_i^2<1,\qquad
 \Psi(\rho,v)=-\sum_i\log(v_i-\rho_i^2)-\log(1-\sum_i v_i^2),
\]

并对v取内部最小值。该域在ρ上的投影精确等于‖ρ‖₄\textless1。每个−log(v−ρ²)的二阶Hessian正定：其对角元为2/d+4ρ²/d²、1/d²，非对角元为−2ρ/d²，行列式为2/d³\textgreater0，其中d=v−ρ²。故Ψ关于(ρ,v)严格凸。

固定内部ρ时，v的有界域边界上至少一个对数发散，故存在唯一内部极小v。记δ=1−Σv\_i²，其站立性给

\[
 2v_i(v_i-\rho_i^2)=\delta,\qquad
 v_i=\frac{\rho_i^2+\sqrt{\rho_i^4+2\delta}}2.
\]

δ+Σv\_i(δ)²=1的左端严格递增，从δ=0时小于1增长到大于1，因而δ唯一。这正是实现消去的辅助变量，未删去原密度方向。

对不同ρ\_1、ρ\_2及其唯一极小v\_1、v\_2，使用联合严格凸性，再在中间ρ处取v最小值，得到约化密度障碍严格凸。ρ逼近L4边界时，δ≤1−Σρ\_i⁴→0，因此约化障碍也发散。

有限μ\textgreater0、正障碍权重、严格可行的有界IR截面上，完整障碍加线性目标取得唯一中心：两点密度不同由上述约化密度障碍严格控制；密度相同而自由变量不同，则至少一个原生Im
f坐标变化，相应盘障碍严格凸。其余凸障碍不破坏严格性，边界发散与有限截面的紧性保证取得性。

这个定理不证明μ→0的极值C唯一，也不把给定Newton容差变成全系数或相移误差界。保存的点只能称达到所列数值标准的中心。

\subsection{4.
恢复、选择与求值的完整C身份}\label{ux6062ux590dux9009ux62e9ux4e0eux6c42ux503cux7684ux5b8cux6574cux8eabux4efd}

原式通过的数值中心逐位保留。若最终转换后的C未通过独立原式核验，仅沿同一IR域内的完整系数线段恢复；任何实际改动都撤销中心资格。固定x的两端属于同一截面，线段仍属于该截面，浮点偏差另报。

\texttt{center\_converged.npz}保存达到中心标准时的全部C；选择接口重新检查该C与JSON/candidate/中心逐项相同、准备数据哈希一致、当前C在xref截面、当前C及协向量的原式支持间隙达标。新selection再记录系数文件哈希；evaluate和相移消费接口核对同一哈希。由此保证新增IR对照不是后验换点得到的曲线。

旧三色IR数据的径向收缩及旧标签范围见\href{../IR_CENTER_IDENTITY_REVIEW_ZH.md}{身份审计}，原文件保留。有限支持的局部几何距离应在交付点实际x坐标解释，不把其微小坐标漂移等同于严格xref中心。

\subsection{5.
已运行的独立检查}\label{ux5df2ux8fd0ux884cux7684ux72ecux7acbux68c0ux67e5}

\href{../verification_IR_slice.json}{验证记录}包含完整来源哈希。106项检查仍在三个测试文件中；新增控制覆盖实际附加解析行上的坐标恒等式、原式渐近对偶与primal、完整中心身份。M3/L2真实CLI经历支持、解析审计、固定xref中心、选择及求值链条；固定截面先在μ=.001求中心再降μ，42步完成。一个直接从μ=1e−6开始而停滞的运行保留，并未被称为不可行。

此处证明的是明确有限IR程序的数学自洽与实现验收原则。3123个采样盘及五个领先条件不是全能区、无限自旋幺正认证。

\newpage

\section{给定S和FF时，有限电流谱存在性的完整判据}\label{ux7ed9ux5b9asux548cffux65f6ux6709ux9650ux7535ux6d41ux8c31ux5b58ux5728ux6027ux7684ux5b8cux6574ux5224ux636e}

范围：原论文每个电流通道的两条有限FESR，采用同一正权重支持集；S、FF及归一化因子固定。以下证明不新增生产约束，也不将谱构造作为相移代表。它补全Schur下界与谱预算之间的逻辑。

\subsection{1.
PSD给出的最小谱}\label{psdux7ed9ux51faux7684ux6700ux5c0fux8c31}

在一个电流节点，写原实Gram为

\[
G=\begin{pmatrix}A&v\\v^T&R\end{pmatrix},\qquad
A=\begin{pmatrix}1+\operatorname{Re}S&\operatorname{Im}S\\
\operatorname{Im}S&1-\operatorname{Re}S\end{pmatrix},\qquad
v=\sqrt2\binom{\operatorname{Re}\mathcal F}{\operatorname{Im}\mathcal F}.
\]

若Δ=1−\textbar S\textbar²\textgreater0，则A正定，Schur补给出必要且充分条件

\[
R\ge r_*=v^TA^{-1}v
=|\mathcal F|^2+\frac{|\mathcal F-S\mathcal F^*|^2}{\Delta}.
\]

这不只是一条较弱的两pion下界。相位失配的完整代价必须保留。

\subsection{2.
两条矩和误差盒的充要判据}\label{ux4e24ux6761ux77e9ux548cux8befux5deeux76d2ux7684ux5145ux8981ux5224ux636e}

令两个矩权重为w\_j\textgreater0及w\_jσ\_j，其中σ\_j\textgreater0。标量通道对应原(0,1)矩，向量通道对应原(−1,0)矩。\textbf{对保存的浮点程序应取σ\_j=W\_\{1j\}/W\_\{0j\}；不能把未证明逐位相等的比值直接替换为节点s\_j。}
理论精确求积中这个比值就是s\_j。

定义最小谱的两矩

\[
a_0=\sum_jw_jr_{*,j},\qquad a_1=\sum_jw_j\sigma_jr_{*,j},
\quad \sigma_-=\min_j\sigma_j,\quad\sigma_+=\max_j\sigma_j.
\]

FESR允许m\_i∈{[}L\_i,U\_i{]}，区间已经包含原始绝对误差。则存在非负剩余谱R\_j−r\_\{*,j\}≥0，使两条矩同时落在误差盒中，\textbf{当且仅当}

\[
\begin{aligned}
a_0&\le U_0, &a_1&\le U_1,\\
\sigma_+a_0-a_1&\le\sigma_+U_0-L_1,
&a_1-\sigma_-a_0&\le U_1-\sigma_-L_0.
\end{aligned}
\]

证明：令t\_i=m\_i−a\_i。剩余谱产生的矩锥恰是

\[
t_0\ge0,\qquad \sigma_-t_0\le t_1\le\sigma_+t_0.
\]

必要性由正权重平均直接成立。充分性在σ\_-\textless σ\_+时将t\_0按比例分配到达到两个端点比值的节点，即得到任意允许t\_1；t\_0=0时必须t\_1=0。σ\_-=σ\_+时锥退化为一条射线，取单个节点即可。

进一步，对误差盒存在可行的m\_0，当且仅当

\[
\max\!\left(L_0,a_0,a_0+\frac{L_1-a_1}{\sigma_+}\right)
\le
\min\!\left(U_0,a_0+\frac{U_1-a_1}{\sigma_-}\right).
\]

将两侧逐项比较，除L\_i≤U\_i外，得到前述四个不等式。剩余的一项
(σ\_+−σ\_-)a\_1≤σ\_+U\_1−σ\_-L\_1
已由a\_1≤U\_1及L\_1≤U\_1推出，因此没有漏掉独立条件。

\subsection{3.
该投影保持凸性}\label{ux8be5ux6295ux5f71ux4fddux6301ux51f8ux6027}

矩阵分式vᵀA⁻¹v在A正定时关于(A,v)联合凸，因为其epigraph正是上述PSD
Gram。S和FF到(A,v)的映射是仿射的。四个左端分别为r\_*的以下非负权重和：

\[
\sum w_jr_{*,j},\quad\sum w_j\sigma_jr_{*,j},\quad
\sum w_j(\sigma_+-\sigma_j)r_{*,j},\quad
\sum w_j(\sigma_j-\sigma_-)r_{*,j}.
\]

所以它们是凸函数，四个次水平集合的交仍凸。若将来使用这项谱消元来改进求解，它可以保留原有限问题，而不需要降低振幅维数、加入ρ目标或丢掉FESR下界。目前生产求解仍显式使用原Gram和谱变量；没有仅凭这份推导宣布新的算法已验证。

\subsection{4.
边界、构造与物理范围}\label{ux8fb9ux754cux6784ux9020ux4e0eux7269ux7406ux8303ux56f4}

在\textbar S\textbar=1处，必须另有v∈Range(A)，再用vᵀA⁺v。否则不存在有限R，不能用无限谱或任意PSD零特征值掩盖。高能无矩节点同样需要此相容条件；假设所有相关节点严格\textbar S\textbar\textless1时，才可无条件作有限Schur提升。

端点节点上的剩余谱分配是\textbf{存在性证明}，不是独立预选的物理电流谱形状。不能把该稀疏构造的尖峰当作预测的ρ，也不能替代原中心选择程序。

连续版本需要另外知道同一S/FF在整个区间上的幺正性、r\_*可积性，并计算真正积分而不是中点求和。非负剩余\textbf{测度}的矩锥可用区间端点得到类似判据；若要求普通密度，端点原子及严格内点情形还需区分。当前有限节点通过没有建立这些连续前提。

因此，完整有限谱可行见证、相位失配的谱代价、以及``所有可行S都出现正确ρ''是不同命题。前两者可用这里的充要判据分析；后者仍须有额外证明或实测比较，不能由两个矩自动推出。

\newpage

\section{从闭式振幅证明全能区IR非空，以及UV约束的非平凡性}\label{ux4eceux95edux5f0fux632fux5e45ux8bc1ux660eux5168ux80fdux533airux975eux7a7aux4ee5ux53cauvux7ea6ux675fux7684ux975eux5e73ux51e1ux6027}

这是一个辅助定理，\textbf{不替代Fig.7的绿色边界代表，也不计作B--F的数值复现}。它回答底层函数族是否连IR约束本身都不自洽，并说明归一化FF、高能界和谱矩为什么会排除足够弱的无ρ振幅。结论关于全部分波的弹性块收缩，不构造其它散射通道的完整解析幺正S矩阵。

\subsection{1.
一个属于当前解析族的明确函数}\label{ux4e00ux4e2aux5c5eux4e8eux5f53ux524dux89e3ux6790ux65cfux7684ux660eux786eux51fdux6570}

取

\[
h(s)=1+z(s)=\frac4{2+\sqrt{4-s}},
\] \[
A_\lambda(s,t,u)=\lambda\left[
\frac{h(s)}8+\frac{h(t)+h(u)}8+
h(s)(h(t)+h(u))+\frac12h(t)h(u)\right],\qquad\lambda>0.
\]

对任意M，当前完整cardinal族都包含此函数：令v\_j=sinφ\_j，T0=0，σ₁=σ₂=λv/8，R=λvvᵀ，Q=λvvᵀ/2。完整DST的第一谐波恒等式给∑v\_jh\_j=h；Q非对角的原存储系数是2Q。这里使用明确的闭式函数，不从有限H证书转移连续结论。

h在规定割平面解析、满足实解析性；A在t↔u下对称，按原同位旋交叉组合得到完整pion振幅。对t\textless4有正Stieltjes表示

\[
h(t)=\int_4^\infty\frac{d\mu(x)}{x-t},\qquad
d\mu(x)=\frac{4\sqrt{x-4}}{\pi x}\,dx.
\]

令x=4+u²，积分化为(8/π)∫u²/{[}(u²+4)(u²+4−t){]}du；部分分式积分即得上述h。它给出所需正性及绝对可积控制，随后由解析延拓固定物理边界值。

\subsection{2.
对任意分波统一控制交叉投影}\label{ux5bf9ux4efbux610fux5206ux6ce2ux7edfux4e00ux63a7ux5236ux4ea4ux53c9ux6295ux5f71}

设d=s−4\textgreater0，t=−d(1−x)/2，定义

\[
J_l=\frac14\int_{-1}^1P_l(x)h(t)dx,
\qquad U_l=\frac14\int_{-1}^1P_l(x)h(t)h(u)dx.
\]

正Stieltjes表示给J\_l=∫dμ(a)
Q\_l(1+2a/d)/d\textgreater0。这里Q\_l(y)\textgreater0可直接由Rodrigues分部积分证明：

\[
Q_l(y)=2^{-l-1}\int_{-1}^1\frac{(1-x^2)^l}{(y-x)^{l+1}}dx>0\quad(y>1).
\]

由\textbar P\_l\textbar≤1，又有0\textless J\_l≤J₀。部分分式

\[
\frac1{(a-t)(b-u)}=\frac{(a-t)^{-1}+(b-u)^{-1}}{a+b+d}
\]

表明奇l时U\_l=0，偶l时0≤U\_l≤2J\_l，因为内部剩余积分是h(−a−d)≤1。交换积分由正测度和\textbar P\_l\textbar≤1的绝对控制保证。

物理割线上

\[
|h(s)|=\frac4{\sqrt s},\quad
\operatorname{Im}h(s)=\frac{4\sqrt{s-4}}s,\quad
J_0=\frac4{s-4}\left[\sqrt s-2-2\log\frac{\sqrt s+2}4\right].
\]

所以J₀≤4/(√s+2)，且J₀≥h(−d)/2=2/(√s+2)。因此

\[
\frac{\operatorname{Im}h}{\kappa J_l}\ge\frac1\pi,
\qquad |h|\le2,\qquad |h|\le4J_0.
\]

\subsection{3.
全部允许分波的幺正不等式}\label{ux5168ux90e8ux5141ux8bb8ux5206ux6ce2ux7684ux5e7aux6b63ux4e0dux7b49ux5f0f}

直接使用三种交叉组合可得（I=0,2取偶l，I=1取奇l）

\[
f_l^0=\lambda\left[\frac5{16}\delta_{l0}h+
(\frac54+9h)J_l+\frac72U_l\right],
\] \[
f_l^2=\lambda\left[\frac18\delta_{l0}h+
(\frac12+3h)J_l+2U_l\right],
\qquad f_l^1=\lambda hJ_l.
\]

从上节界限，统一有

\[
|f_l^I|\le\lambda C_IJ_l,
\qquad \operatorname{Im}f_l^I\ge\lambda c_I\operatorname{Im}h\,J_l,
\quad (C_0,C_1,C_2)=(55/2,2,11),\quad(c_0,c_1,c_2)=(9,1,3).
\]

S波中额外的直接h项，其虚部为正；实模界用\textbar h\textbar≤4J₀吸收。于是

\[
1-|S_l^I|^2\ge
\kappa\lambda J_l\left(2c_I\operatorname{Im}h-\kappa\lambda C_I^2J_l\right)>0
\]

只需

\[
0<\lambda<\min_I\frac{2c_I}{\pi C_I^2}
=\frac{72}{3025\pi}.
\]

同一个λ因此对所有s\textgreater4、所有允许l有效；不是逐个l取一个不同的小参数。阈值、s→∞或l→∞时裕量可以趋零，未声称存在统一正裕量。

同时\textbar f\_l\textsuperscript{1\textbar≤λ，故\textbar S\_l}1−1\textbar≤πλ\textless1。P1的S始终在右半平面，连续相移没有90°上穿。这是明确的无ρ式相移族，而不是用图像挑出的无峰点。

\subsection{4.
同时满足给定IR范数}\label{ux540cux65f6ux6ee1ux8db3ux7ed9ux5b9airux8303ux6570}

在四个手征输入点及其角积分中，s,t,u∈{[}0,4)，所以1≤h\textless2，从而\textbar A\_λ\textbar≤43λ/4。利用同位旋系数及∫\textbar P₁\textbar=1，有

\[
|f_0^0|\le215\lambda/8,\quad
|f_1^1|\le43\lambda/8,\quad |f_0^2|\le43\lambda/4.
\]

四点上\textbar R₀₁\textsuperscript{χ\textbar≤9/2、\textbar R₂₁}χ\textbar≤9/7，故两条独立L2残差分别不超过817λ/8和989λ/28。取λ\textless8εχ/817即可使两球严格满足。实际双密度又满足‖r‖₄≤λ{[}M²+M(M+1)/2{]}\^{}(1/4)，因此再取λ小于B除以此因子即可。

这证明当前解析函数族在原IR不等式层面有明确的全能区、全允许分波例子；不是完整QCD或完整多通道S矩阵的存在定理，也没有得到物理fπ位置的边界代表。

\subsection{5.
原UV输入为何排除该族的弱耦合部分}\label{ux539fuvux8f93ux5165ux4e3aux4f55ux6392ux9664ux8be5ux65cfux7684ux5f31ux8026ux5408ux90e8ux5206}

只用向量通道。记F\_j=1+∑K\_ji v\_i+i
v\_j，v\_i=ImF\_i，\(\mathcal F_j=k_jF_j\)。低能集合L取s\_j≤s0，高能集合H取s\_j\textgreater s0；向量高能上界给\textbar v\_j\textbar、\textbar ReF\_j\textbar≤c\_j=√ε\_F/k\_j（ε\_F=3×10⁻⁵）。

对任意h∈H，归一化与三角不等式给

\[
\left|\sum_{j\in L}K_{hj}v_j\right|\ge
b_h:=1-c_h-\sum_{j\in H}|K_{hj}|c_j.
\]

若b\_h\textgreater0，这要求低能ImF不能全小。另一方面，本族有1−ReS\_j≤πλ。原Gram相应2×2主子式给

\[
R_j\ge\frac{2k_j^2v_j^2}{1-\operatorname{Re}S_j}
\ge\frac{2k_j^2v_j^2}{\pi\lambda}.
\]

向量逆矩上界U及正权重w\_j由此限制∑\_L w\_j
k\_j²v\_j²≤πλU/2。加权Cauchy--Schwarz于是得到必要条件

\[
\lambda\ge\frac{2b_h^2}{\pi U D_h},\qquad
D_h=\sum_{j\in L}\frac{K_{hj}^2}{w_jk_j^2}.
\]

对原M50、打印矩/raw绝对.002/hard-midpoint、s0=(1.2/.14)²及向量FF上界，使用精确原生节点、K和
k\_j²=(s\_j−4)√(1−4/s\_j)/(384π⁵)、w\_j=(π/M)tan(φ\_j/2)，逐项区间计算给第49号（零基）高节点的

\[
b_h>0.85464142575,\qquad
\lambda\ge2.35634303968\times10^{-5}.
\]

而λ=2⁻³⁰满足上述全部严格IR条件，却违反此UV必要下界。因此该明确完整解析IR振幅没有原有限UV扩展；实际上整个上述保证手征可行的小λ区间都被这个界排除。全部高节点、公式输入与包络见\href{../GLOBAL_IR_UV_EXCLUSION.json}{计算记录}。没有输入任何ρ质量或论文输出曲线。

\subsection{6.
此定理闭合了什么}\label{ux6b64ux5b9aux7406ux95edux5408ux4e86ux4ec0ux4e48}

它严格连接了``共同解析、交叉及分波收缩的IR族存在''\,``IR不强制ρ式上穿''和``归一化FF＋高能FF界＋有限谱预算会排除其中的弱相互作用部分''。它没有证明所有无ρ振幅都被UV排除，更没有从一个逆矩推出唯一ρ质量。

实际Fig.7绿色边界振幅的全部C另有原式及解析Farkas排除证据；完整多变量UV代表、图形比较和分辨率检查也另行交付。这个辅助解析例子不替代那些主线结果。

\newpage

\section{不借助输出曲线的UV→P波必要下界}\label{ux4e0dux501fux52a9ux8f93ux51faux66f2ux7ebfux7684uvpux6ce2ux5fc5ux8981ux4e0bux754c}

本节把\href{../../../evidence/GLOBAL_IR_SEED_AND_UV_OBSTRUCTION_20260911.md}{明确IR族的UV排除}中的关键步骤推广到任意原有限UV可行点。只使用原向量电流Gram、F(0)=1所固定的离散色散关系、向量高能FF界和逆矩上界；不使用代表坐标、实验ρ质量、相移或正则化数值拟合。

记低能节点集L=\{s\_j≤s0\}，高能集H=\{s\_j\textgreater s0\}，F\_j=1+Σ\_i
K\_ji v\_i+i
v\_j，v\_i=ImF\_i，ℱ\_j=k\_jF\_j，R\_j为向量电流谱。记w\_j\textgreater0为低能逆矩权重，Σ\_L
w\_jR\_j≤U。高能界给\textbar F\_j\textbar≤c\_j=√(3×10⁻⁵)/k\_j。

对每个h∈H，

\[
\left|\sum_{j\in L}K_{hj}v_j\right|
\ge b_h:=1-c_h-\sum_{j\in H}|K_{hj}|c_j.
\]

取b\_h\textgreater0。完整Gram的一个2×2主子式给

\[
2k_j^2v_j^2\le R_j y_j,\qquad
 y_j=1-\Re S_j=\kappa_j\Im f_j,\quad0\le y_j\le2.
\]

这条不等式在y\_j=0或R\_j=0也成立；此时相关v\_j=0，无需除以零。加权Cauchy--Schwarz于是严格推出

\[
 b_h^2\le
 \left(\sum_{L}w_jR_j\right)
 \left(\sum_L\frac{K_{hj}^2 y_j}{2w_jk_j^2}\right)
 \le\frac U2\sum_L\frac{K_{hj}^2y_j}{w_jk_j^2}.
\]

因此任何可行P1振幅都满足一个显式线性必要条件

\[
\boxed{\sum_L\frac{K_{hj}^2\kappa_j\Im f_j}{w_jk_j^2}
\ge\frac{2b_h^2}{U}}.
\]

令D\_h=Σ\_L K\_hj²/(w\_jk\_j²)，则至少有一个低能节点满足

\[
\max_{j\in L}(1-\Re S_j)
\ge\frac{2b_h^2}{U D_h}>0.
\]

先前\href{../GLOBAL_IR_UV_EXCLUSION.json}{精确K／节点的区间记录}已经给出原M50、printed/raw/hard条件下全部高节点的b\_h、D\_h和U。对h=49，右端为该记录lambda\_lower乘π，严格大于7.40×10⁻⁵。因此至少一个原生低能P1节点有

\[
\frac{|S_j-1|^2}{4}\ge\frac{(1-\Re S_j)^2}{4}
>1.369\times10^{-9}.
\]

这里的``吸收''指振幅的虚部κImf，不是非弹性损失1−\textbar S\textbar²。上述结果是很弱但严格、且不依赖任何边界选择的吸收／散射强度下界。它说明UV输入为何不能由无限减弱整个P波来满足：FF归一化和高能衰减要求低能ImF承担有限色散贡献，而Gram及有限正谱预算把这转成P波吸收的需求。

这个下界仍远小于90°上穿所需的1−ReS\textgreater1（S≠0时）。所以\textbf{该不等式自身不证明ρ共振、峰位、峰数或极点}；也不能据此声称完整UV可行域包含或不包含无ρ振幅。论文的几何预选边界、实际相移曲线和五组分辨率检验仍是另一层必须完成的数值结果。

它适用于原有限UV模型及其加强子集。原生K、k、权重和Γ的身份必须相同；不同插值、矩单位或密度定义之间不能转移该数值下界。

\newpage

\section{有限正则化的充分结论与实测非平凡依赖}\label{ux6709ux9650ux6b63ux5219ux5316ux7684ux5145ux5206ux7ed3ux8bbaux4e0eux5b9eux6d4bux975eux5e73ux51e1ux4f9dux8d56}

本文件补充``加入一个有限B即可称正则化已处理完毕''的不足。讨论的是固定M、L、H、全部IR/UV输入与规范之后的有限联合问题；不改变原基线，不用新B挑选相移。

\subsection{一、定义、有限性与曲率}\label{ux4e00ux5b9aux4e49ux6709ux9650ux6027ux4e0eux66f2ux7387}

令X=(C,ImF,R\_current)为完整变量，r(C)是实际双密度的独立项（对称非对角储存系数先除2）。除了密度界以外的全部约束固定，定义

\[
V(B)=\sup\{f_0^0(3):X\text{满足这些原约束},\ \|r(C)\|_4\le B\}.
\]

实际双密度界结合原生S0/S2的自由列恒等式控制单密度及T0，从而控制全部C；固定有限B时目标有有限上界。这不是由双密度界单独推出，也不要求整个H全列满秩。完整电流变量可能无界，所以这里用sup，不擅自假设最优完整电流谱取得。

V在非空可行域上非减。它也凹：取B1、B2下两个完整可行点X1、X2，凸组合Xθ=θX1+(1−θ)X2仍满足散射盘、手征球、PSD
Gram、FF界及FESR；F(0)=1作为仿射归一化亦保持。三角不等式给

\[
\|r(C_\theta)\|_4\le\theta B_1+(1-\theta)B_2.
\]

目标线性，故对任意逼近两个sup的可行序列取极限，得到

\[
V(\theta B_1+(1-\theta)B_2)\ge\theta V(B_1)+(1-\theta)V(B_2).
\]

这里必须混合完整X，不能把已归一化FF的联合问题当作齐次问题而只缩放C。有限凹函数在其非空定义域内部连续且局部Lipschitz；这并不保证B→∞的有限极限，更不证明相移或ρ读数对B不敏感。

\subsection{二、固定对偶给出整段B的上界}\label{ux4e8cux56faux5b9aux5bf9ux5076ux7ed9ux51faux6574ux6bb5bux7684ux4e0aux754c}

原式支持核验先消去未受直接范数控制的单密度及T0残差。对剩余实际双密度协向量g，Hölder不等式给

\[
g\cdot r\le B\|g\|_{4/3}.
\]

将一个已经保存的对偶固定，其余圆盘、手征、电流谱、FF及矩支持项合记为a。它们在这里都不依赖B：电流残差界来自固定FESR和FF输入；Gram合同尺度也是固定协向量定义的一部分。因此对所有B′≥0，均有

\[
V(B')\le a+B'\|g\|_{4/3}.
\]

若原报告的上界为U\_B、密度支持项为D\_B=B‖g‖\_(4/3)，同一协向量给

\[
V(B')\le U_B+(B'/B-1)D_B.
\]

不要求该协向量在新B′下最优。实际数值结论仍须用原式区间运算重验，避免把打印小数的差值冒充向外舍入证书。

\subsection{三、原M50/L10模型的直接重验}\label{ux4e09ux539fm50l10ux6a21ux578bux7684ux76f4ux63a5ux91cdux9a8c}

基线B=377500。原+x tip有

\[
V(B)\ge0.09903980867034552,
\qquad U_B=0.09911418047407121,
\qquad D_B=0.0111350014628817.
\]

采用B′=B/2=188750，只重放已有tip协向量，并用已有D3完整联合见证检验新可行域；没有新优化或选点。新原式重验得到

\[
0.04085344847737422\le V(B/2)\le0.09354667974263037.
\]

D3的实际双密度范数约174150.23158644，低于188750；全部原式联合约束通过。随后包含解析行误差的独立primal也通过1500散射盘、100个Gram、2条χ、4条FESR、14个FF界和L4。它证明减半B后仍有真实的有限解析族节点见证。

因而同一有限H程序中严格有

\[
V(B)-V(B/2)\ge0.0054931289277>0.
\]

这不是``点恰好靠近范数边界''的经验判断：它比较的是原B的已知可行下界与减半B后所有可行点的严格上界。原tip的双密度范数约376307.43也接近B，但这个接近本身不是敏感性证明，上述两侧证据才是。

证据：\href{../regularizer_half_bound_replay/report.json}{减半B的原式重验}、\href{../regularizer_half_analytic/source_audit.json}{解析primal}、\href{../../sequential_reproduction_20260910/E1_tip_center/report.json}{原tip}。机器可核验的差值及来源哈希见\href{../REGULARIZATION_VALUE_PROOF.json}{数值汇总}。

\subsection{四、对复现结论的影响}\label{ux56dbux5bf9ux590dux73b0ux7ed3ux8bbaux7684ux5f71ux54cd}

原2309正文§3.1给出采样与变量，§3.2只写``some
norm''；它没有唯一给出本复现中\(B(M)=100[M^2+M(M+1)/2]\)这一配方。因此不能说这套正则化已证明无关紧要，也不能在缺少作者完整数值规范时期待相同投影边界由已列物理参数唯一决定。\href{https://arxiv.org/abs/2309.12402v3}{原论文}

这项审计支持把正则化规范视为重要的复现输入缺口；它没有证明该缺口单独造成了P1质量偏差、Fig.8形状或全部L依赖。当前没有改变B，没有运行B扫描，没有用B′生成新相移或匹配论文输出。

严格目标上界在本文件中属于已保存的有限H/K/FESR程序。解析源可行点另通过核验；没有把有限上界转移成全能量、全自旋或连续物理振幅的证书。条件性的有限紧性、凹性与灵敏度结论不能替代这些不同层面的物理验收。

\newpage

\section{ρ式信号：峰存在的严格弱结论及其限度}\label{ux3c1ux5f0fux4fe1ux53f7ux5cf0ux5b58ux5728ux7684ux4e25ux683cux5f31ux7ed3ux8bbaux53caux5176ux9650ux5ea6}

本文件只使用已经在相移读取之前冻结的原生UV三代表，没有新优化或代表选择。对其P1散射强度及向量FF平方的\textbf{全部原生局部极大候选}，用保存的解析f包络、精确原生κ和精确cardinal
Hilbert矩阵重新核验了中点高于两侧端点的严格间隙。

令J(E)为P1强度κ²\textbar f₁¹\textbar²/4或\textbar F₁\textbar²。若a\textless b\textless c且已证明

\[
J(b)>\max(J(a),J(c)),
\]

由同一有限解析函数的连续性，J在{[}a,c{]}取得最大值，且最大值不能在端点，故(a,c)中至少有一个内点峰。这比把折线顶点直接叫作真实峰更强，但它不确定峰的精确位置、唯一性或极点。

本次所有候选都通过严格包络比较：tip的P1有1个，mid/ref各有2个；向量FF平方分别有4、15、15个。原显示能区内，最大原生样本的中点能量仍是tip约792.14MeV、mid/ref约680.41MeV；实际峰只被上述邻点区间包围，不能把这个中点或折线插值当作精确峰质量。

全部候选（包括低能FF小起伏与约1.06GeV的次级P1起伏）都保留在\href{../RHO_SIGNAL_PEAK_EXISTENCE.json}{区间证据}，没有只留下接近期望ρ质量的峰。额外的FF小峰不能自动解释为新介子。\href{../../sequential_reproduction_20260910/CE_native_comparison/rho_mechanism.pdf}{IR/UV同一冻结振幅比较图}展示了原生有限程序中的信号变化。

这些严格弱结论\textbf{不等于正确物理ρ已复现}。该三振幅的离节点幺正违例和无穷远障碍仍然成立。电流R目前是离散变量，其折线峰不能借用这里对解析S/FF的连续性证明。90°读数也仍是原生括区与描述性插值，未升级为一个经过区间验证的负实S轴交叉或第二张面极点证书。

由此可以精确回答``ρ是否出现''：已定义解析振幅确有P1强度和FF的内点峰，原生相移也显示90°括区；但质量、形状、稳健性和连续物理有效性仍是分别待验的更强主张。不能因Fig.7的IR对照没有合适ρ而判失败，也不能因UV中有峰就宣布全部核心主张完成。

\newpage

\section{同模型IR→UV的严格分离：复用已有对偶，不另作优化}\label{ux540cux6a21ux578biruvux7684ux4e25ux683cux5206ux79bbux590dux7528ux5df2ux6709ux5bf9ux5076ux4e0dux53e6ux4f5cux4f18ux5316}

\subsection{对象和结果}\label{ux5bf9ux8c61ux548cux7ed3ux679c}

IR点是在物理参考xref=.07332139057293466处按最大化f11(3)预选的完整3876系数振幅。其数值中心、原式支持和解析primal均通过，且交付C与保存中心逐位相同。它与UV程序共用同一个3123盘准备、T0=0、五个领先高能必要条件、两条separate-L2手征球(.002)和实际双密度L4(B=377500)。只有UV加入归一化FF、电流Gram、打印FESR/raw绝对.002及原高能FF界；本份证书还明确采用F0、F1在z=−1处至少二阶零的端点完成，即两者的值与一阶导数均为零。这些FF端点等式不是原2309有限样本界自动包含的约束，不能删除后直接搬用本负界。

\href{../interlaced_IR_ref_02/report.json}{IR支持}、\href{../interlaced_IR_ref_analytic/source_audit.json}{IR解析核验}、\href{../interlaced_IR_selected/selection.json}{相位前选择}。这里尚不借助该IR点的相移是否有ρ来选点或作证明。

复用已完成的\href{../interlaced_UV_ref_support_01/report.json}{UV参考截面支持}的电流对偶，对这个IR点作零目标Farkas复核，得到

\[
U_{\rm fixed,H}=-0.1850389032871912<0.
\]

进一步计入原生解析行包络、精确运动学及FF/FESR输入，得到

\[
\boxed{U_{\rm fixed,analytic}=-0.18503890326560488<0.}
\]

后一个数是向外舍入的严格上界，完整包络和来源见\href{../interlaced_fixed_IR_analytic/report.json}{解析Farkas记录}。两次复核均没有新优化。这是\textbf{该完整IR振幅没有满足所声明UV条件的电流扩展}的证书，不是完整联合域不可行------同一域已有其他完整可行见证。

\subsection{为什么可以复用UV支持的对偶}\label{ux4e3aux4ec0ux4e48ux53efux4ee5ux590dux7528uvux652fux6301ux7684ux5bf9ux5076}

固定C后，电流Gram具有形式

\[
G_q(C,v,R)=G_{q,0}+G_{q,C}(C)+G_{q,v}(v)+G_{q,R}(R)\succeq0.
\]

对任意已审计的半正定对偶Y\_q，Tr(Y\_qG\_q)≥0。把这些不等式与FESR区间的支持函数、FF圆盘的支持函数相加，并对尚未精确消去的电流变量残差使用独立上界，可得

\[
0\le c_0+a\cdot C+r_v\cdot v+r_R\cdot R
\le c_0+a\cdot C+h_{\mathcal B_v}(r_v)+h_{\mathcal B_R}(r_R)
=:U_{\rm fixed}(C).
\]

这些残差上界来自正矩权重、Gram所给R≥\textbar ℱ\textbar²、高能FF界和FF归一化；自由高能谱没有虚构上界，其残差符号必须合法。FF端点依赖变量按原仿射关系精确消去。必要时对保存的浮点Y作向外PSD修复并重新计算整个界，不能仅信任求解器状态。

原UV线性支持证书使用的正是这些电流协向量，并再加入散射、手征和L4的支持函数以消去自由C。固定C后不必重新寻找协向量；把其余C变量消元改为直接代入，就得到一个有效的零目标上界。等价地，原UV支持平面若严格排除该IR点，其差距可以分解为上述固定C上界与已满足的散射/手征/密度约束的非负支持余量，因此已有电流对偶已经提供一个可用的分离方向。

最终验收仍依靠完整的重新计算。本例得到U\_fixed\textless0；若真的存在电流扩展，则必须同时满足0≤U\_fixed\textless0，矛盾。对偶的负数大小可以通过整体缩放改变，不应解释为物理距离或ρ质量误差。

\subsection{原式与解析复核的区别}\label{ux539fux5f0fux4e0eux89e3ux6790ux590dux6838ux7684ux533aux522b}

第一份证书使用保存的float64
H。第二份在同一个完整解析函数族中包住固定C的原生分波值，并用精确K、运动学、打印十进制矩和hard-midpoint权重复核电流证书；没有把旧H证书直接搬成解析结论。

作为复核时的任意电流候选，程序可读取UV点的v、R；这不声称它们对IR点可行。零目标上界对所有允许v、R成立，结论不依赖这个任意候选。IR选择目录没有加入joint.npz或current.json，因而后续IR图仍明确只有散射/手征输入。

\href{../fixed_IR_dual_before_fix/report.json}{接口修复前的失败}保留：旧dual入口错误要求IR目录也有joint.npz。新入口用单独的UV对偶供体并固定所给IR的C，正确执行已有数学证书；\href{../interlaced_fixed_IR_dual/report.json}{原式重验}和解析重验分别保存。随后补丁加入工作进程读取前的物理算子输入哈希。新解析复核有完整算子哈希；较早的原式供体复核仍保留其inputs为空的历史记录，不能追认成当时已经记录了读取前哈希。源码快照及新解析证书的完整输入保持可复核。

\subsection{不能由此推出的结论}\label{ux4e0dux80fdux7531ux6b64ux63a8ux51faux7684ux7ed3ux8bba}

这只排除一个预定IR振幅；不排除整个IR集合，不证明所有无ρ振幅都不可能，也不证明UV允许域只有一个ρ极点。ρ是否在所选UV点形成正确而稳定的结构，仍由同一完整振幅的相移、强度、FF、电流谱及分辨率检验回答。

\subsection{删除新增FF端点条件后的独立复核}\label{ux5220ux9664ux65b0ux589effux7aefux70b9ux6761ux4ef6ux540eux7684ux72ecux7acbux590dux6838}

为避免把上述加强模型证书搬到较弱模型，又单独复用了保留的原基线E2\_ref\_verified对偶，固定同一个新IR完整C，在原A2\_M50\_L10/D1\_M50\_L10准备上重新计算。原1500散射行及十一辅助行与加强准备的原生部分逐位相同，连解析误差包络也相同，见\href{../ORIGINAL_NATIVE_PREFIX_REPLAY.json}{直接重验}；没有跨PV／analytic处方转移证据。

这次电流模型不含新增的FF端点零条件，也不含新增高能散射约束；F(0)=1、原电流Gram、打印/raw/hard的四矩及原高能FF样本界保留。原式上界为−.1484359517640356，进一步\href{../original_UV_fixed_new_IR_analytic/report.json}{计入解析误差}后仍有

\[
\boxed{U_{\rm original\ UV,analytic}=-0.1484359517403425<0.}
\]

因此原有限UV条件已足以排除这个新IR振幅，不依赖新增FF二阶端点条件。这里是另一份完整重算的证书，不是把较强条件下的负上界直接沿用；两份对偶不同，其负数大小也不用于度量条件强弱。它仍只排除该C，不证明所有无ρ振幅都被排除。

\newpage

\section{电流Gram为什么允许90°上穿而没有正确的近弹性ρ}\label{ux7535ux6d41gramux4e3aux4ec0ux4e48ux5141ux8bb890ux4e0aux7a7fux800cux6ca1ux6709ux6b63ux786eux7684ux8fd1ux5f39ux6027ux3c1}

这里从原Gram直接推导，不增加η≈1、谱饱和或目标质量约束。下面的真实例子已通过3582盘修补之前的3123盘有限模型的全部原式与解析条件；该模型仍有独立5MeV网格违例，不能称为完整连续QCD反例。

\subsection{1.
原Gram的等价圆盘}\label{ux539fgramux7684ux7b49ux4ef7ux5706ux76d8}

在一个物理节点，设R\textgreater0，ℱ≠0，并记

\[
r=\frac{|\mathcal F|^2}{R},\qquad Q=\frac{\mathcal F}{\mathcal F^*},\qquad |Q|=1.
\]

原复Gram的两个电流主子式给0≤r≤1，散射主子式给\textbar S\textbar≤1。行列式条件等价于

\[
r|S-Q|^2\le(1-r)(1-|S|^2).
\]

展开并配方得到

\[
\boxed{|S-rQ|\le1-r.}
\]

反过来，这个圆盘与0≤r≤1给\textbar S\textbar≤1及全部原复Gram主子式非负，因此与Gram半正定等价。\textbar S\textbar=1的边界由连续性或直接的核兼容条件覆盖。F=0时r=0、rQ可直接写作ℱ²/R=0，恢复单位圆盘；R=0时PSD要求F=0。

这说明当前耦合约束允许的S，是以rQ为中心、1−r为半径的圆盘。它严格推出

\[
\eta=|S|\ge\max(0,2r-1),
\]

却没有一般地强制η≈1。只有r≈1时圆盘才缩到接近弹性点Q。r是电流谱的两pion份额，不是散射的弹性概率η²，也不是某个共振的分支比。

\subsection{2.
一个明确的局部反例，以及实际完整点}\label{ux4e00ux4e2aux660eux786eux7684ux5c40ux90e8ux53cdux4f8bux4ee5ux53caux5b9eux9645ux5b8cux6574ux70b9}

取R=1、ℱ=i√(3/5)、S=−1/5。于是r=3/5、Q=−1，S恰在允许圆盘的内侧边界。实Gram为

\[
G=\begin{pmatrix}
4/5&0&0\\
0&6/5&\sqrt{6/5}\\
0&\sqrt{6/5}&1
\end{pmatrix}\succeq0.
\]

它的特征值为0、4/5、11/5，秩二；相移为90°，但η=.2，散射强度\textbar S−1\textbar²/4=.36。故一个零特征值、单个90°相位读数或Watson半相位吻合（arg
F=arg
S/2，即S与Q同相），都不足以推出近弹性ρ。单点本身不证明上穿；若取局部光滑递增的δ(s)，并令S(s)=.2
exp(2iδ(s))、ℱ(s)=√.6
exp(iδ(s))、R=1，则在δ穿过90°的局部区间上这个Gram仍处处半正定。这仍是局部代数例子，不是完整QCD构造。秩一/两pion饱和、秩二边界和弹性S=Q必须分开。

本次三份完整预选UV振幅也实测到这种行为。下表取各自预先冻结C的最大原生P1强度样本；节点并不用于重新选C：

{\def\LTcaptype{none} % do not increment counter
\begin{longtable}[]{@{}lrrrrr@{}}
\toprule\noalign{}
点 & E(GeV) & η & r & Gram允许的最小η & 半相位差 arg(S/Q)/2 \\
\midrule\noalign{}
\endhead
\bottomrule\noalign{}
\endlastfoot
tip & .731675 & .304105 & .591962 & .183925 & −19.9252° \\
mid & .680414 & .214759 & .606548 & .213095 & −2.86759° \\
ref & .680414 & .196052 & .598010 & .196020 & −.402572° \\
\end{longtable}
}

\href{../GRAM_DISK_DIAGNOSTIC.json}{区间核算}重新使用保存的解析f包络、精确K/k和FF端点完成，三个圆盘裕量均严格正。尤其ref几乎贴着允许圆盘的内侧边界，完全符合原式，不是幺正约束漏掉、相位分支错误或一次Newton未收敛造成的。

当前三点的90°读数约752/680/679MeV只是原生插值指标；5MeV网格读数约747/680/680MeV。它们有显著FF和谱增强，但质量分散、强非弹性和S0定量偏差仍是major问题，不能只挑tip较接近770MeV就宣称成功。

\subsection{3. 哪一层还欠缺}\label{ux54eaux4e00ux5c42ux8fd8ux6b20ux7f3a}

F(0)=1、高能FF界和有限矩预算严格排除某些IR振幅，也能迫使非零P波吸收。然而它们不自动把这个圆盘压到外侧弹性点。当前raw逆矩误差约47\%，且只给两个谱矩；即使矩精确，也仅约束加权平均而非全部谱形。完整数值例子说明``PSD近边界''不能代替实际的谱份额和η检查；这并未反证原文的特定数值曲线。

在固定Q时，Re(Q*
S)在上述圆盘上的唯一最大点是S=Q，最大值为1。因此用旧FF生成Q、在保留全部约束的条件下最大化此类线性泛函，是一种明确的弹性偏好选择原则。全幅度的交叉/解析约束可能阻止各节点同时达到这个局部理想点；Q随新FF更新后也不能仅凭一次最优证明得到固定点或唯一性。

作者\href{https://arxiv.org/html/2403.10772v4}{2403后续工作的迭代目标}提供了这种具名方向，本仓库已有对应CLI。它属于后续的选点算法，不是2309原有限条件本身强制近弹性的证明。任何实际尝试都须保留全部原约束，记录坐标变化、旧Q目标、完整新C、原式支持和新S/新F残差；不能加入目标ρ质量或将相位替换成FF相位来掩盖η问题。

\newpage

\section{暂缓问题：有限解析族不能精确表达``先弹性、后非弹性''的阈值切换}\label{ux6682ux7f13ux95eeux9898ux6709ux9650ux89e3ux6790ux65cfux4e0dux80fdux7cbeux786eux8868ux8fbeux5148ux5f39ux6027ux540eux975eux5f39ux6027ux7684ux9608ux503cux5207ux6362}

本条是更强连续物理目标的适用范围证明，\textbf{不作为2309有限节点程序的新增验收前提，也不据此修改生产参数或约束}。它不反证论文的有限精度图形结果。

\subsection{命题及证明}\label{ux547dux9898ux53caux8bc1ux660e}

在当前固定有限M的函数族中，每个基函数是

\[
z(s)=\frac{2-\sqrt{4-s}}{2+\sqrt{4-s}}
\]

的有限多项式，振幅由单变量项、两变量乘积以及常数的有限和组成。物理分波通过固定区间x∈{[}−1,1{]}上的Legendre投影得到，其中t=−(s−4)(1−x)/2，u=−(s−4)(1+x)/2。

\textbf{命题。} 对任意固定有限系数C、任意固定分波l，函数

\[
\Delta_l(s)=1-|1+i\pi\sqrt{1-4/s}\,f_l(s)|^2
\]

在整个实区间(4,∞)上实解析。因此，若它在任意非空开子区间上恒等于零，则它在全部(4,∞)上恒等于零。

证明：任取s\emph{\textgreater4。物理边界值中的√(s−4)可在s}的足够小复邻域选取解析分支。对所有x∈{[}−1,1{]}，t、u在实s\emph{处均不大于零，故4−t、4−u与其平方根奇点具有一致正距离；缩小s}邻域后，这个距离仍对全角区间一致成立。所有有限基函数及其乘积遂在该邻域解析，并在紧角区间上一致有界。固定有限Legendre投影保留解析性。运动学因子在s*\textgreater4附近也解析，故ReS、ImS及其平方和在实轴上实解析。实解析函数在连通区间上的恒等定理给出后半结论。证毕。

这里只使用物理边界值的局部解析延拓，不宣称原第一张面跨越整条右割线无割。

\subsection{物理含义与不能外推的结论}\label{ux7269ux7406ux542bux4e49ux4e0eux4e0dux80fdux5916ux63a8ux7684ux7ed3ux8bba}

如果更强的目标要求某一分波在最低多粒子阈值以下\textbf{精确}弹性，而在更高能量有严格非零的粒子产生损失，那么任何固定有限M的当前解析族都不能同时精确实现这两点。真实多粒子阈值处的非解析性在这一个有限多项式族中没有单独出现。

这不是``所有非零有限振幅都不可能''的证明：全程弹性的特殊函数不由本命题排除，单纯满足\textbar S\textbar≤1的非平凡有限振幅也已明确构造。命题更不否认用有限基在有限精度下近似物理振幅，或通过M→∞的非一致解析极限恢复阈值结构。当前bootstrap只施加必要的收缩盘，没有把一个开放能段上的\textbar S\textbar=1作为精确等式。

在理想同位旋pion模型中常讨论4π阈值s=16；本证明不依赖这个具体数，只需要目标存在一个真正的弹性开区间和后续非弹性区间。是否允许额外更轻通道属于另一项物理输入，不能从给出的有限盘自行推出。

\subsection{可行的后续路线}\label{ux53efux884cux7684ux540eux7eedux8defux7ebf}

若目标升级为精确连续的多通道QCD振幅，函数族应显式包含相应的多粒子割线，或对含物理阈值的连续谱作可控极限；还要处理弹性等式和非弹性通道之间的关系。增加物理阈值变量、分段色散谱或多通道/N-over-D参数化都是可研究方向，但需要新的来源、收敛与交叉验证，不能只把当前有限点证书改名。

本轮仍按论文有限原型推进：修复已经测到的离节点违例，保存严格有限primal/dual、原参数比较、相移和分辨率证据。上述更强问题另行暂缓，不用来替代或拖延这些明确的科学计算。

\end{document}
